\section{Methodology}

The implemented pipeline consists of five stages: data preparation, factorized model training, Rashomon-set selection, prediction-disagreement analysis, and explanation-variability analysis. All stages are exposed through the \texttt{mmcyber} command line interface.

\subsection{Data Preparation}

The current implementation uses NSL-KDD as the default intrusion-detection dataset. The dataset is downloaded automatically on the first run and stored locally under \texttt{data/raw/nsl-kdd/}. The task is configured as binary classification by default, mapping all attack labels to one attack class and retaining normal traffic as the second class. A multiclass configuration is supported in the data loader but is not the focus of the current draft.

Categorical features, namely protocol type, service, and flag, are transformed with one-hot encoding. Numeric features are standardized with a standard scaler. The preprocessing pipeline is fitted on the training split and then applied to validation and test data. The fitted preprocessor and feature names are stored in the run directory to keep later explanation steps aligned with the trained models.

\subsection{Model Families}

The primary model family is a feed-forward multilayer perceptron implemented in PyTorch. The architecture is configurable through the experiment configuration file. Hidden layers use linear transformations, a configurable activation function, batch normalization, and dropout, followed by a final linear classification layer. MLPs are trained with AdamW and cross-entropy loss. Early stopping is based on validation loss.

In addition, the pipeline includes XGBoost classifiers as a separate comparison family. These models are not used to decompose within-family multiplicity factors such as activation functions or hidden-layer widths. Instead, they provide an architecture-level comparison against the MLP-based Rashomon analyses.

\subsection{Factorized Candidate Model Set}

To induce and separate sources of model multiplicity, the pipeline uses a factorized design centered on a fixed \texttt{origin\_seed}. This anchor seed defines the run identity, the baseline MLP, and the naming of all artifacts written to disk. As a result, repeated experiments with different origin seeds can be grouped cleanly for later statistical testing.

Starting from the baseline MLP, the implementation generates additional candidate models in four comparison blocks. First, a seed block varies the training seed while keeping the data, architecture, activation, and training fraction fixed. Second, a bootstrap block keeps the training seed fixed at the origin seed but changes the resampled training set through bootstrap resampling. Third, a joint architecture--activation block crosses several hidden-layer configurations with several activation functions. The baseline combination is excluded from this block because it is already represented by the reference model. Fourth, a separate XGBoost block trains multiple boosted-tree models with identical hyperparameters but different random seeds.

This design intentionally separates within-MLP multiplicity from between-family comparisons. The seed and bootstrap blocks isolate stochastic effects, the joint architecture--activation block captures deliberate modeling choices within the MLP family, and the XGBoost block provides an external point of comparison rather than another factor inside the MLP decomposition.

Each trained model is evaluated on the same held-out test split, and the pipeline stores per-model metrics as well as per-sample predictions and predicted class probabilities. In addition, the explicit training plan is exported so that every model can be traced back to its originating block, origin seed, training seed, data seed, model family, architecture, and activation setting.

\subsection{Training-Data Perturbations}

The factorized design distinguishes two types of data variation. Ordinary training runs use stratified subset sampling so that class balance is approximately preserved. In contrast, the bootstrap block samples with replacement from the training set to create alternative empirical datasets while keeping the remaining modeling decisions fixed. This distinction allows the analysis to separate variability due to stochastic optimization from variability due to perturbations of the observed training sample.

\subsection{Rashomon-Set Selection}

After training, models are filtered into a Rashomon set. Given a metric such as accuracy and a tolerance value, the best observed score is identified and all models with score at least the best score minus the tolerance are retained. The tolerance therefore controls how broadly similarly good models are included in the analysis.

\subsection{Predictive Multiplicity}

Predictive multiplicity is measured at two levels. At the model-pair level, the pipeline computes pairwise disagreement rates between all models in the Rashomon set. Each model pair is also labeled by the factor block from which it originates whenever both models belong to the same non-baseline comparison block. This yields block-level summaries such as seed, bootstrap, architecture--activation, and XGBoost disagreement.

At the sample level, predictions are pivoted into a model-by-sample vote matrix. For each sample, the implementation computes the majority prediction, vote entropy, the number of unique predictions, the majority fraction, and the conflict ratio. The conflict ratio is defined as one minus the majority fraction and is positive whenever at least one Rashomon model disagrees with the majority vote.

\subsection{Explanation Variability}

The pipeline currently computes SHAP explanations for each MLP model in the Rashomon set. A random subset of training samples is used as background data, and a configurable subset of test samples is explained. When enabled, the \texttt{--only-conflicts} option restricts the explanation set to samples with positive conflict ratio, matching the conflict-focused analysis used in the prior BA-style experiments.

Although SHAP is model-agnostic in principle and can also be applied to tree-based models such as XGBoost, the current implementation excludes XGBoost artifacts from the explanation stage. This is therefore a scope decision of the present pipeline rather than a limitation of the SHAP framework itself. The restriction keeps the explanation analysis homogeneous while the predictive multiplicity comparison across model families is established first.

For each model, class, sample, and feature, the SHAP value is stored in long format. The implementation also computes mean absolute SHAP values per feature and model, which are used to compare feature rankings and explanation distances. Pairwise explanation disagreement is summarized using cosine distance between mean-absolute-SHAP vectors and Jaccard overlap between top-\(k\) feature sets.

\subsection{Variability Metrics}

The BA-style analysis aggregates SHAP values across models for each sample, class, and feature. It computes the minimum, maximum, mean, standard deviation, variance, and range of SHAP values. In addition, sign instability is measured as the smaller of the positive and negative sign fractions across models. This captures whether similarly accurate models agree on the direction of a feature's contribution.

Finally, the implementation relates explanation variability to predictive multiplicity by computing Spearman correlations between conflict ratio and SHAP value range, SHAP value variance, and sign instability for each feature and class. The plotting module produces publication-oriented figures for metric variability, predictive disagreement, SHAP feature importance, explanation disagreement, sign-instability heatmaps, correlation matrices, and per-feature scatter plots.
